% Part: first-order-logic
% Chapter: beyond
% Section: modal-logics

\documentclass[../../../include/open-logic-section]{subfiles}

\begin{document}

\olfileid[tr]{fol}{byd}{mod}

\olsection{Kiplik Mantıkları}

Şu koşullu tümce örneğini ele alın:
\begin{quote}
  Jeremy o odada yalnızsa sarhoş ve çıplaktır ve sandalyelerin üzerinde dans
  ediyordur.
\end{quote}
Bu, maddi bakımdan doğru olabilen ama yine de yanıltıcı olan bir koşullu sav
örneğidir; çünkü öncülle sonuç arasında yalnızca öncülün yanlış ya da sonucun
doğru olmasından daha güçlü bir bağ bulunduğunu düşündürür. Başka bir deyişle
ifade, savın yalnızca bu belirli dünyada (Jeremy odada yalnız olmadığı için
önemsiz biçimde doğru olabileceği dünyada) doğru olduğunu değil, dahası öncül
doğru \emph{olsaydı} sonucun da doğru \emph{olacağını} düşündürür. Yani savın,
yalnızca bu dünyada değil her ``olanaklı'' dünyada doğru olduğu; ya da yalnızca
bu belirli dünyada doğru olmanın tersine \emph{zorunlu olarak} doğru olduğu
anlamına alınabilir.

Kiplik mantığı bu tür zorunluluğu anlamlandırmak üzere tasarlanmıştır. Olağan
önermeler mantığına bir kutu işleci eklenerek modal önermeler mantığı elde
edilir; yani $!A$ !!a{formula} ise $\Box !A$ da öyledir. Sezgisel olarak
$\Box !A$, $!A$'nın \emph{zorunlu olarak} doğru ya da her olanaklı dünyada
doğru olduğunu savlar. $\Diamond !A$ genellikle $\lnot \Box \lnot !A$'nın
kısaltması sayılır ve $!A$'nın doğru olmasının \emph{olanaklı} olduğunu
savladığı biçiminde okunabilir. Elbette kiplik, yüklem mantığına da eklenebilir.

Kripke !!{structure}{}sı kiplik mantığına bir semantik sağlamak için
kullanılabilir; aslında Kripke bu semantiği ilk kez kiplik mantığını düşünerek
tasarlamıştır. Kısmi sıralamalarla sınırlanmak yerine, daha genel olarak bir
``olanaklı dünyalar'' kümesi $P$ ve dünyalar arasında ikili bir
``erişilebilirlik'' bağıntısı $\Atom{R}{x,y}$ vardır. Sezgisel olarak
$\Atom{R}{p,q}$, $q$~dünyasının $p$~dünyasıyla bağdaşabilir olduğunu savlar;
yani $p$~dünyasında ``bulunuyorsak'', dünyanın $q$~dünyası gibi olabileceği
olasılığını hesaba katmamız gerekir.

Kiplik mantığına, ``kaplamsal'' mantığın karşıtı olarak bazen ``içlemsel''
mantık denir. Klasik mantık gibi kaplamsal bir mantığın amaçlanan semantiği,
yalnızca tek bir dünyaya, ``edimsel'' dünyaya gönderimde bulunur; ``içlemsel''
bir mantığın semantiğiyse daha ayrıntılı bir ontolojiye dayanır. Zorunluluğu
!!{structure}{}landırmanın yanı sıra kiplik, uygulamaya göre $\Box$ ile $\Diamond$
yeniden yorumlanarak başka dilsel kuruluşları da !!{structure}{}landırmak için
kullanılabilir. Örneğin:
\begin{enumerate}
\item Kanıtlanabilirlik mantığında $\Box !A$, ``$!A$ kanıtlanabilir'' ve
  $\Diamond !A$, ``$!A$ tutarlıdır'' diye okunur.
\item Epistemik mantıkta $\Box !A$, ``$!A$'yı biliyorum'' ya da ``$!A$'ya
  inanıyorum'' diye okunabilir.
\item Zamansal mantıkta $\Box !A$, ``$!A$ her zaman doğrudur'' ve
  $\Diamond !A$, ``$!A$ bazen doğrudur'' diye okunabilir.
\end{enumerate}

Mantığı, kipliği ele alan kurallar ve aksiyomlarla genişletmek isteriz. Örneğin
$\Log{S4}$ sistemi, önermeler mantığının olağan aksiyomları ve kurallarıyla
birlikte şu aksiyomlardan oluşur:
\begin{align*}
& \Box (!A \lif !B) \lif (\Box !A \lif \Box !B)\\
& \Box !A \lif !A \\
& \Box !A \lif \Box \Box !A
\intertext{ve ayrıca ``$!A$'dan $\Box !A$ sonucunu çıkar'' kuralından oluşur.
  $\Log{S5}$ ise şu aksiyomu ekler:}
& \Diamond !A \lif \Box \Diamond !A
\end{align*}
Bu aksiyomların çeşitlemeleri farklı uygulamalara uygun olabilir; örneğin S5'in
genellikle mantıksal zorunluluk kavramını nitelediği kabul edilir. Güzel olan
şudur ki !!{derivation} sisteminin sağlam ve tam olduğu bir semantik, Kripke
!!{structure}{}sındaki erişilebilirlik bağıntısını doğal biçimlerde kısıtlayarak
genellikle bulunabilir. Örneğin $\Log{S4}$, erişilebilirlik bağıntısı
yansımalı ve geçişli olan Kripke !!{structure}s sınıfına karşılık gelir.
$\Log{S5}$ ise erişilebilirlik bağıntısı {\em evrensel}, yani her dünyanın
öteki her dünyadan erişilebilir olduğu Kripke !!{structure}s sınıfına karşılık
gelir; dolayısıyla $\Box !A$ ancak ve ancak $!A$ her dünyada geçerliyse
geçerlidir.

\end{document}
